% ---
% id: meat-core/frank-wolfe/req:bracket-on-reference
% kind: req
% title: "The first-order bound evaluated on the reference companion"
% summary: >-
%   Where a reference companion is admitted, the run's first-order number is
%   the best bracket on the companion at the best points of the snapped search
%   and of a short search continued on the companion from them, each assessed
%   by the usual gate; the snapped bound is recorded as a diagnostic only.
% prominence: standard
% status: implemented
% uses: [cor:fo-reference, thm:firstorder, cor:fo-keylength, req:first-order-withdrawal, req:iterate-checks, req:support-rank-check, req:acceptance-geometries, req:acceptance-iterate-exact, req:verified-feasible-side, req:first-order-clamp, req:bound-estimate-difference, meat-core/req:companion-admission, meat-core/prop:reference-companion, foundations/def:evaluated-off-snap, foundations/req:snap-reference]
% mentions: [rem:fo-arithmetic, req:certified-evaluation, req:negative-estimate-refused, meat-core/prop:optol-nobound, foundations/ass:snapped-measurement]
% assumes: []
% refs: []
% provenance: newly-derived
% implements: certiqs_key.solvers.frank_wolfe.reference:bracket_on_reference
% lean:
% notes: "The second cause of (iv), factor_enclosure_inconsistent, withdraws the bracket rather than only its certification because the first-order number reads the entrywise enclosure, which an inconsistency puts in doubt; the snapped bound reads neither enclosure. The polish's share of the wall budget is set aside before the search on D begins, so that a search stopped by its budget, where the snapped statistics also slow it, still leaves the polish its time, at the cost of a shorter search on D; its points enter the certified evaluation through (i) as the search's do, which is why it forms its acceptance-side iterate on the exact island."
% ---

The search runs on the snapped declaration; the bound is formed on the
companion. Any points at which Corollary~\ref{cor:fo-reference} holds give a
valid bound, so the choice of points bears on its size only. The companion
and the snapped declaration differ only in the measured families of the
factors the companion replaces and in the entries of their last links, so the
points the snapped bracket ranks highest are the natural candidates, and a
few of them keep the extra oracle calls few; the last point is added because
it is the point the search ended at. Those points can still sit badly on the
companion. The search fits the acceptance-side component to the snapped
statistics, and where the snap moved an entry that is small beside its
outcome's probability, the companion's statistics at the same point differ
from that component by a relative amount of the move's size, which the
factor $\alpha/(\alpha-1)$ of the objective magnifies: the bracket on the
companion can then fall far below the snapped one, though the two values are
close. The polish, a short search continued on the companion's own program
from the best of those points, moves the point to where the companion's
bracket is large. Its points are points like any other, since
Corollary~\ref{cor:fo-reference} reads of a point nothing but its
hypotheses, and the gate decides them at each point whatever produced it:
in exact arithmetic the points it adds can only raise a maximum of valid
brackets. That is not the whole of it at the first-order tier. Where the
companion carries a positive radius, its first-order bracket is formed from
the binary64 numbers nearest the centres of its enclosure, which are
displaced from the companion itself by an amount no claim bounds
(Remark~\ref{rem:fo-arithmetic}, Proposition~\ref{prop:optol-nobound}); a
search that maximises that evaluation makes it more likely that such an
unbudgeted excess reaches the first-order number than points found on
another declaration do. The polish's points count at the first-order tier
all the same, as the search's do; where every radius is zero only the
rounding of the evaluation is exposed, as at any point; and the certified
tier, formed from the exact centres and radii
(Requirement~\ref{req:certified-evaluation}), is unaffected. Which blocks then rest
on no premise about the snap follows from what the companion carries, block
by block (\xref{foundations/def:evaluated-off-snap}). It follows also from
which number the key length is formed from: where the companion's bound
counts but its certified evaluation certifies no point, a number certified
on the snapped declaration may carry the key length instead
(Requirement~\ref{req:certified-evaluation}), and the blocks the companion
carried then rest on the snap premise after all.

\begin{requirement}\label{req:bracket-on-reference}
Where the engine has admitted a compiled reference companion
$\Decl^{\mathrm{ref}}$ beside the declaration $\Decl$ it searched on
(Requirement~\ref{req:companion-admission}), it forms the run's first-order
number as follows, after the search on $\Decl$ has ended.
\begin{enumerate}[label=(\roman*),leftmargin=2.4em]
\item \emph{The points.} It evaluates the points of two sets. The first
  holds at most five points of the search on $\Decl$: the at most four
  members of the index set of $\Decl$'s bound
  (Requirement~\ref{req:first-order-withdrawal}) with the largest brackets
  on $\Decl$, ties broken by the order of the run, and the last point the
  search assessed; where $\Decl$'s index set is empty, the last point alone.
  The second holds the points kept by the \emph{polish}, a search on the
  Choi program of $\Decl^{\mathrm{ref}}$ that continues the search on
  $\Decl$ after it has ended. The polish reads the objective, the gradient,
  the maps and the entries of $\Decl^{\mathrm{ref}}$
  (Requirement~\ref{req:companion-admission}\,(i)) and no map generated from
  $\Decl$'s measured family; it starts at the point of the first set with
  the largest bracket on $\Decl$, at the last point where $\Decl$'s index
  set is empty, and forms its acceptance-side iterate as
  Requirement~\ref{req:acceptance-iterate-exact} requires, from that point's
  exact acceptance-side component. Where that component was not formed on
  the \tref{exact island} (Requirement~\ref{req:acceptance-iterate-exact}),
  the polish does not run, the second set is empty and the record names
  that reason. It assesses each point it consumes by the
  gate of (ii) and keeps the at most four of them that contribute with the
  largest brackets on $\Decl^{\mathrm{ref}}$, ties broken by its own order,
  and the last point it assessed; a polish that assesses no point keeps
  none. It runs within the share \iface{reference_polish_share} of the
  wall budget of one evaluation of the run, where the run has one, set
  aside for it before the search on $\Decl$ begins, and ends where that
  share is spent, where an iteration cap it records is reached or where its
  own stopping rule holds, whichever comes first; a polish that fails or is
  cut short keeps, as above, the points it assessed until then. The share
  so set aside shortens the search on $\Decl$, which bears on the run's
  estimate, on $\hlb[\Decl]$, on a number certified on $\Decl$ in the
  fallback of Requirement~\ref{req:certified-evaluation}\,(i), and on the
  soundness of none of them. Its objective values are not the run's
  \tref{estimate}: none is reported, and none refuses the run
  (Requirement~\ref{req:negative-estimate-refused} reads the estimate of the
  search on $\Decl$ alone).
  Each point $x_t=(\mathbf q_t,J_t)$ of either set is taken as the search
  or the polish consumed it, $\mathbf q_t$ included, and is a point of the
  Choi program of $\Decl^{\mathrm{ref}}$
  (Corollary~\ref{cor:fo-reference}\,(ii)).
\item \emph{The gate.} At each such point it forms the bracket of
  \eqref{eq:hlb} for $\Decl^{\mathrm{ref}}$: the objective of
  $\Decl^{\mathrm{ref}}$ at $x_t$, its gradient at $x_t$ with no
  perturbation, the acceptance-side minimisation of
  Requirements~\ref{req:acceptance-geometries}
  and~\ref{req:acceptance-iterate-exact}, and the feasible side of
  Requirement~\ref{req:verified-feasible-side}, from a fresh feasible-side
  oracle call on the linear form of that gradient over
  $\Decl^{\mathrm{ref}}$'s entries. It decides (H-POS) and (H-SUPP) on
  $\Decl^{\mathrm{ref}}$'s maps (Requirement~\ref{req:support-rank-check})
  and makes every check of Requirement~\ref{req:iterate-checks} against
  $\Decl^{\mathrm{ref}}$'s entries, on the point; a point at which
  Requirement~\ref{req:first-order-withdrawal} withdraws it on these terms
  contributes no bracket, with that requirement's reason.
\item \emph{The number.} Where some evaluated point contributes and the
  bracket is not withdrawn by (iv), the run's
  first-order number is $\hlb[\Decl^{\mathrm{ref}}]$, the largest bracket so
  formed, its clamp and key length are those of
  Requirement~\ref{req:first-order-clamp} and
  Corollary~\ref{cor:fo-keylength} at it, with the parameters of $\Decl$'s
  run, and its tier is \emph{first-order}. The bound on $\Decl$,
  $\hlb[\Decl]$, is recorded beside it with its brackets as a diagnostic,
  and no key length is formed from it. The difference
  $\hlb[\Decl^{\mathrm{ref}}]-\hlb[\Decl]$ is recorded with its sign, as a
  diagnostic, beside the difference of
  Requirement~\ref{req:bound-estimate-difference}, which is taken against
  the estimate of $\Decl$. A point of the second set that contributes
  counts as any evaluated point does: where no point of the first set
  contributes, it makes this item apply in place of the fallback of (iv),
  so that the run's first-order number moves from $\hlb[\Decl]$ to the
  companion.
\item \emph{The bracket withdrawn.} The bracket on the companion is
  withdrawn for one of two causes: no evaluated point contributes; or the
  certified evaluation finds the two enclosures of a block the companion
  carries inconsistent, the cause \iface{factor_enclosure_inconsistent}
  (Requirement~\ref{req:certified-evaluation}\,(iv)). An inconsistency shows
  that a module carrying the enclosures is defective, and the first-order
  number on the companion reads its entrywise enclosure, which may be the
  defective one; the second cause is decided after that number was formed,
  and the run then forms its first-order number again as this item says,
  before any key length is reported. Where the bracket is withdrawn, the
  run's first-order number is $\hlb[\Decl]$, formed as without a
  companion, with the reason \iface{reference_bracket_withdrawn}, its cause
  and the reasons of every evaluated point; no block is then evaluated off its snap.
  The record names in the group \emph{reference bracket withdrawn}, with the
  cause (\xref{foundations/req:snap-reference}\,(iii)(b)), every moved block
  on which a declaring instance whose link the companion replaces handed the
  reference and every moved class of each snapped emission of such a link
  whose normalizer the companion bypasses; it names as it does without a
  companion every other moved block or class that rests on
  \xref{foundations/ass:snapped-measurement}, so that a moved block on which a declaring
  instance whose link the companion does not replace handed a reference
  rests in the group \emph{reference bracket not evaluated}
  (\xref{foundations/req:snap-reference}\,(iii)(c)). A run
  that declares the strict option \iface{strict_snaps} refuses there instead
  and reports no key length (\xref{foundations/req:snap-reference}\,(iv)).
\item \emph{The blocks.} Where (iii) applies and the key length the run
  reports is not formed in the fallback of
  Requirement~\ref{req:certified-evaluation}\,(i), the record of the run names as
  \tref{evaluated off its snap} (\xref{foundations/def:evaluated-off-snap})
  every \tref{moved} block of the snap of each declaring instance whose
  link the companion replaces, on which the companion carries the handed
  reference, and every data sector of each snapped emission of such a link
  that holds a moved class and whose normalizer the companion bypasses; a kept block or sector is named as kept and is not
  counted among them. It names every moved block on which the companion
  carries the snapped block as not so evaluated, so that it rests on
  \xref{foundations/ass:snapped-measurement}
  (\xref{foundations/req:snap-reference}\,(iii)), and a run that declares
  \iface{strict_snaps} refuses there (\xref{foundations/req:snap-reference}\,(iv)).
  Where (iii) applies but the key length the run reports is formed from a
  number certified on $\Decl$ in the fallback of
  Requirement~\ref{req:certified-evaluation}\,(i), no number formed on the
  companion reaches that key length, and the number it is formed from reads
  the snapped blocks: the moved blocks on which the companion carries the
  handed reference and the data sectors whose normalizer it bypasses are
  then not evaluated off their snap, condition (ii) of
  \xref{foundations/def:evaluated-off-snap} failing for them. The record
  names them as resting on \xref{foundations/ass:snapped-measurement} in
  the group \emph{reference bracket withdrawn} with the cause
  \iface{companion_not_certified}
  (\xref{foundations/req:snap-reference}\,(iii)(b)), beside the moved blocks
  on which the companion carries the snapped block, which rest as above. The
  run decides this after the certified evaluation and before any key length
  is reported, and no block is then named as evaluated off its
  snap. A run that declares \iface{strict_snaps} forms no such
  number (Requirement~\ref{req:certified-evaluation}\,(i)).
\item \emph{The record.} The record carries, per evaluated point, the set
  of (i) it belongs to, its bracket on $\Decl^{\mathrm{ref}}$ and the
  withdrawal reasons there, and for a point of the first set its bracket on
  $\Decl$; for the polish, whether it ran and, where it did not, why, the
  point it started from, the share of the wall budget set aside for it and
  the wall time it took, its iteration cap, the number of points it
  assessed and the cause it stopped on; the measurements of
  Requirement~\ref{req:companion-admission}\,(iii), the largest radius of the
  companion's enclosure among them, which bear on the arithmetic of
  Remark~\ref{rem:fo-arithmetic} and on nothing else; and, per block,
  whether the companion
  carries the handed reference or the snapped block.
\end{enumerate}
A run with no admitted companion forms its first-order number on $\Decl$
alone, as the other requirements of this submodule state; a moved block on
which a reference was handed, and a moved class of an emission beside which a
companion that bypasses its normalizer was formed, then rest in the group
\emph{reference bracket not evaluated}
(\xref{foundations/req:snap-reference}\,(iii)).
\end{requirement}
